% GENERATED FILE. Edit canonical lessons, then run npm run generate:books.
% canonical-source-hash: fnv1a64:56bff0347cc489f3
% canonical-lessons: TE-C299-sanagalu, TE-C299-avalu, TE-C299-kottimira, TE-C299-karivepaku, TE-C299-majjiga

\chapter{Chickpeas, Mustard seeds, Coriander leaves, Curry leaves, Buttermilk}
\label{ch:te-chickpeas-mustard-seeds-coriander-leaves-curry-leaves-buttermilk}
\hlchaptermodality{299}

\begin{chapteropening}
\textbf{By the end of this chapter:} \emph{I can say chickpeas, mustard seeds, coriander leaves, curry leaves and buttermilk.}

Complete the last lesson of chapter 299: I can say chickpeas, mustard seeds, coriander leaves, curry leaves and buttermilk.
\end{chapteropening}

\section[\texorpdfstring{\te{శనగలు} (śanagalu)}{శనగలు (śanagalu)}]{\te{శనగలు} (śanagalu) — chickpeas}

\label{lesson:TE-C299-sanagalu}



\begin{quote}
Before the new one: say the Telugu for hollow, then the Telugu for dim.
\end{quote}

\subsection*{You'll want to know: \te{శనగలు}}
\textbf{\te{శనగలు}} — \emph{śanagalu} — \textquotedblleft{}chickpeas\textquotedblright{}.

\begin{grammarlens}[title={What you've built}]
One more word for this chapter.
\end{grammarlens}

\subsection*{Your turn}
\begin{itemize}
\raggedright
  \item \emph{Say it:} \emph{śanagalu}
  \item \emph{Say it:} \emph{śanagalu}, once more
  \item \emph{Say it:} say \emph{śanagalu}
  \item \emph{Recall:} say the Telugu for hollow, then the Telugu for dim, then say \emph{śanagalu} again
\end{itemize}

\begin{tcolorbox}[breakable,colback=teal!4,colframe=teal!35!black,title={Before you move on}]
What does \te{శనగలు} mean? (\textquotedblleft{}Chickpeas\textquotedblright{}.) Say it once more.
\end{tcolorbox}
\section[\texorpdfstring{\te{ఆవాలు} (āvālu)}{ఆవాలు (āvālu)}]{\te{ఆవాలు} (āvālu) — mustard seeds}

\label{lesson:TE-C299-avalu}



\begin{quote}
Before the new one: say the Telugu for dim, then the Telugu for chickpeas.
\end{quote}

\subsection*{You'll want to know: \te{ఆవాలు}}
\textbf{\te{ఆవాలు}} — \emph{āvālu} — \textquotedblleft{}mustard seeds\textquotedblright{}.

\begin{grammarlens}[title={What you've built}]
One more word for this chapter.
\end{grammarlens}

\subsection*{Your turn}
\begin{itemize}
\raggedright
  \item \emph{Say it:} \emph{āvālu}
  \item \emph{Say it:} \emph{āvālu}, once more
  \item \emph{Say it:} say \emph{āvālu}
  \item \emph{Recall:} say the Telugu for dim, then the Telugu for chickpeas, then say \emph{āvālu} again
\end{itemize}

\begin{tcolorbox}[breakable,colback=teal!4,colframe=teal!35!black,title={Before you move on}]
What does \te{ఆవాలు} mean? (\textquotedblleft{}Mustard seeds\textquotedblright{}.) Say it once more.
\end{tcolorbox}
\section[\texorpdfstring{\te{కొత్తిమీర} (kottimīra)}{కొత్తిమీర (kottimīra)}]{\te{కొత్తిమీర} (kottimīra) — coriander leaves}

\label{lesson:TE-C299-kottimira}



\begin{quote}
Before the new one: say the Telugu for chickpeas, then the Telugu for mustard seeds.
\end{quote}

\subsection*{You'll want to know: \te{కొత్తిమీర}}
\textbf{\te{కొత్తిమీర}} — \emph{kottimīra} — \textquotedblleft{}coriander leaves\textquotedblright{}.

\begin{grammarlens}[title={What you've built}]
One more word for this chapter.
\end{grammarlens}

\subsection*{Your turn}
\begin{itemize}
\raggedright
  \item \emph{Say it:} \emph{kottimīra}
  \item \emph{Say it:} \emph{kottimīra}, once more
  \item \emph{Say it:} say \emph{kottimīra}
  \item \emph{Recall:} say the Telugu for chickpeas, then the Telugu for mustard seeds, then say \emph{kottimīra} again
\end{itemize}

\begin{tcolorbox}[breakable,colback=teal!4,colframe=teal!35!black,title={Before you move on}]
What does \te{కొత్తిమీర} mean? (\textquotedblleft{}Coriander leaves\textquotedblright{}.) Say it once more.
\end{tcolorbox}
\section[\texorpdfstring{\te{కరివేపాకు} (karivēpāku)}{కరివేపాకు (karivēpāku)}]{\te{కరివేపాకు} (karivēpāku) — curry leaves}

\label{lesson:TE-C299-karivepaku}



\begin{quote}
Before the new one: say the Telugu for mustard seeds, then the Telugu for coriander leaves.
\end{quote}

\subsection*{You'll want to know: \te{కరివేపాకు}}
\textbf{\te{కరివేపాకు}} — \emph{karivēpāku} — \textquotedblleft{}curry leaves\textquotedblright{}.

\begin{grammarlens}[title={What you've built}]
One more word for this chapter.
\end{grammarlens}

\subsection*{Your turn}
\begin{itemize}
\raggedright
  \item \emph{Say it:} \emph{karivēpāku}
  \item \emph{Say it:} \emph{karivēpāku}, once more
  \item \emph{Say it:} say \emph{karivēpāku}
  \item \emph{Recall:} say the Telugu for mustard seeds, then the Telugu for coriander leaves, then say \emph{karivēpāku} again
\end{itemize}

\begin{tcolorbox}[breakable,colback=teal!4,colframe=teal!35!black,title={Before you move on}]
What does \te{కరివేపాకు} mean? (\textquotedblleft{}Curry leaves\textquotedblright{}.) Say it once more.
\end{tcolorbox}
\section[\texorpdfstring{\te{మజ్జిగ} (majjiga)}{మజ్జిగ (majjiga)}]{\te{మజ్జిగ} (majjiga) — buttermilk}

\label{lesson:TE-C299-majjiga}



\begin{quote}
Before the new one: say the Telugu for coriander leaves, then the Telugu for curry leaves.
\end{quote}

\subsection*{You'll want to know: \te{మజ్జిగ}}
\textbf{\te{మజ్జిగ}} — \emph{majjiga} — \textquotedblleft{}buttermilk\textquotedblright{}.

\begin{grammarlens}[title={What you've built}]
One more word for this chapter.
\end{grammarlens}

\subsection*{Your turn}
\begin{itemize}
\raggedright
  \item \emph{Say it:} \emph{majjiga}
  \item \emph{Say it:} \emph{majjiga}, once more
  \item \emph{Say it:} say \emph{majjiga}
  \item \emph{Recall:} say the Telugu for coriander leaves, then the Telugu for curry leaves, then say \emph{majjiga} again
\end{itemize}

\begin{tcolorbox}[breakable,colback=teal!4,colframe=teal!35!black,title={Before you move on}]
What does \te{మజ్జిగ} mean? (\textquotedblleft{}Buttermilk\textquotedblright{}.) Say it once more.
\end{tcolorbox}
